%%% FIELD rate-definition
Put \(N=n+m\) and \(\ell_n=\log(en)\).  The unequal-information annotation-frontier rate is
\[
r_\epsilon(n,m,d):=\min\!\left\{1,\frac1n+\frac{d^2}{N^2\ell_n^2}\right\}.
\]

%%% FIELD estimator-definition
Fix constants \(H_\epsilon>1\), \(c_\epsilon^\circ>0\), and \(n_\epsilon^\circ\) as calibrated in \(\mathrm{lem:poisson\mbox{-}hybrid\mbox{-}risk}\).  For \(n<n_\epsilon^\circ\), set \(\widehat\tau^{\mathrm{mix}}_{n,m,d}=0\).  Otherwise split the labeled records deterministically into three blocks, use the first for outcome marks, and pool the \((X,A)\) projections of the other two with a deterministic split of the auxiliary records to form independent pilot and factorial-moment blocks.  On each block independently average over Poisson prefix lengths having mean one eighth of the available block length, with output zero on prefix overflow.  For a realized valid prefix, let \(u,t_p,t\) be the three Poisson means, put \(L=\max\{2,\lfloor c_\epsilon^\circ\log(en)\rfloor\}\), \(B=H_\epsilon L/\min(t_p,t)\), and call cell \(x\) light when its pilot count is at most \(\lfloor t_pB/4\rfloor\).  Define
\[
 H_L(z)=\frac{1-T_L(1-2z)}{2L^2},\qquad E_L(z)=\frac{H_L(z)}z,qquad G_L(z)=\frac{1-E_L(z)}z,
\]
using the polynomial continuations at zero.  Expand
\[
 W_{a,L}(s_0,s_1)=\frac{s_0+s_1}{B}G_L(s_a/B)
   =\sum_{r_0,r_1}c_{a,r_0,r_1}s_0^{r_0}s_1^{r_1}.
\]
If \(K_{0x},K_{1x}\) are the factorial-block counts, set
\[
 \widehat W_{a,L,x}=\sum_{r_0,r_1}c_{a,r_0,r_1}
 \frac{(K_{0x})_{r_0}(K_{1x})_{r_1}}{t^{r_0+r_1}},
 \qquad
 \widehat W^H_{a,x}=\frac{K_{0x}+K_{1x}+1}{K_{ax}+1}.
\]
For the outcome-block success count \(S_{ax}\), the cell statistic is
\[
 Z_x=\mathbf 1\{x\text{ light}\}\left(\frac{S_{1x}}u\widehat W_{1,L,x}-\frac{S_{0x}}u\widehat W_{0,L,x}\right)
 +\mathbf 1\{x\text{ heavy}\}\left(\frac{S_{1x}}u\widehat W^H_{1,x}-\frac{S_{0x}}u\widehat W^H_{0,x}\right).
\]
The valid-prefix output is \(\operatorname{clip}_{[-1,1]}(\sum_xZ_x)\).  Finally, \(\widehat\tau^{\mathrm{mix}}_{n,m,d}\) is its conditional expectation over the independent prefix lengths given \((\mathcal L_n,\mathcal U_m)\).  This is a finite weighted sum, hence a total computable measurable map of the observed samples.

%%% FIELD prior-definition
Let \(N=n+m\) and \(\ell_n=\log(en)\).  In the regime \(\min\{1,d^2/(N^2\ell_n^2)\}\le C_\epsilon/n\), let \(\Pi_0,\Pi_1\) be point masses at the two laws concentrated on \(X=1\), with \(e_1=1/2\), \(\mu_{01}=0\), and \(\mu_{11}=1/2\mp h_n\), where \(h_n=c_\epsilon/\sqrt n\le1/4\).  In the complementary regime use the following approximation-dual construction.  Put \(L=\lceil C_\epsilon^\circ\ell_n\rceil\), \(B=b_\epsilon L/N\), \(a=\gamma_\epsilon B/L^2\), and \(\kappa=(1-2\epsilon)/\epsilon\), with the constants calibrated in \(\mathrm{lem:common\mbox{-}marginal\mbox{-}poisson\mbox{-}prior}\).  On \(I=[a/\kappa,B]\), let \(\sigma\) be the finite signed measure of total variation one which annihilates \(1,p,\ldots,p^L\) and for which \(\int p/(p+a)\,d\sigma(p)\ge c_\epsilon>0\).  Write \(h=d\sigma/d|\sigma|\), and define
\[
 \nu(dp)=\frac{a}{p+a}|\sigma|(dp)\quad(p\in I),
 \qquad
 \nu(\{0\})=1-\int_I\frac{a}{p+a}|\sigma|(dp).
\]
Choose \(k=\min\{d-1,\lfloor c_\epsilon^\circ NL\rfloor\}\) rare cells.  Draw their raw masses independently from \(\nu\), put \(s_i=\epsilon(p_i+a)\) when \(p_i>0\), and set \(e_i=s_i/p_i\).  The reservoir raw mass is \(p_*=1-kE_\nu[p]\), with propensity \(1/2\) and zero outcomes.  Divide every raw cell mass by \(S=p_*+\sum_{i=1}^kp_i\).  Under \(\Pi_1\), put \(\mu_{1i}=(1+h(p_i))/2\); under \(\Pi_0\), put \(\mu_{1i}=(1-h(p_i))/2\); under both put \(\mu_{0i}=0\), and assign identical legal values at raw mass zero.  Remaining alphabet cells have mass zero.  In either branch \(\Pi_0\) and \(\Pi_1\) are probability measures on \(\mathcal M_{d,\epsilon}\) and induce exactly the same random \(P_{XA}\).

%%% FIELD inverse-poisson-statement
Let \(u,t>0\).  In the independent Poisson experiment with outcome-success counts \(S_{ax}\sim\operatorname{Poi}(u q_{ax})\) and marginal counts \(K_{ax}\sim\operatorname{Poi}(t s_{ax})\), define
\[
 Z^H=\sum_{x=1}^d\left\{\frac{S_{1x}}u\frac{K_{0x}+K_{1x}+1}{K_{1x}+1}-
 \frac{S_{0x}}u\frac{K_{0x}+K_{1x}+1}{K_{0x}+1}\right\}.
\]
Uniformly over \(P\in\mathcal M_{d,\epsilon}\),
\[
 |E_PZ^H-\tau(P)|\le C_\epsilon\min\{1,d/t\},
 \qquad
 \operatorname{Var}_P(Z^H)\le C_\epsilon(u^{-1}+t^{-1}).
\]

%%% FIELD inverse-poisson-proof
Fix a cell and first take arm \(a=1\).  Write \(s=pe\), \(q=s\mu\), \(K=K_{1x}\), and \(V=K_{0x}\).  Independence of the Poisson coordinates follows from Poisson splitting.  Shifting the absolutely convergent Poisson series gives
\[
 E\frac1{K+1}=e^{-ts}\sum_{j\ge0}\frac{(ts)^j}{j!(j+1)}
 =\frac{1-e^{-ts}}{ts},
\]
with the continuous value one at \(ts=0\).  Therefore, for \(W=1+V/(K+1)\),
\[
 EW=1+t(p-s)\frac{1-e^{-ts}}{ts}
 =\frac1e-\frac{1-e}{e}e^{-tpe},
\]
and, using \(E(S/u)=q\),
\[
 E\{(S/u)W\}=p\mu\{1-(1-e)e^{-tpe}\}.
\]
The occupied-cell overlap assumption gives
\[
 |E Z^H-\tau(P)|
 \le2\sum_xp_xe^{-\epsilon t p_x}.
\]
Since \(z e^{-\epsilon tz}\le(\epsilon et)^{-1}\) for \(z\ge0\), the last sum is at most \(2d/(\epsilon et)\); it is also at most two because \(\sum_xp_x=1\).  This proves the bias assertion.

For the variance, the integral identity \((j+1)^{-1}=\int_0^1z^j\,dz\) and Tonelli's theorem for nonnegative summands imply
\[
 E(K+1)^{-2}=\int_0^1\!\int_0^1e^{-ts(1-zw)}\,dz\,dw
 \le \frac{C}{(1+ts)^2}.
\]
The same argument with two additional integrations gives
\[
 E\{(K+1)(K+2)\}^{-2}\le \frac{C}{(1+ts)^4}.
\]
Indeed, on writing \(z=1-r,w=1-v\), one has \(1-zw\ge(r+v)/2\), which proves the first bound; repeating this comparison for the four integration variables proves the second.  Consequently, using \(EV^2=t(p-s)+t^2(p-s)^2\) and overlap,
\[
 EW^2\le C_\epsilon.
\]

For completeness, the Poisson variance inequality used next follows by exposing the points of each independent Poisson process one at a time: the Doob martingale differences are conditional expectations of the add-one increment, and conditional Jensen followed by summation of their quadratic variations yields
\[
 \operatorname{Var}f(K,V)\le ts\,E(\Delta_Kf)^2+t(p-s)\,E(\Delta_Vf)^2.
\]
Here
\[
 \Delta_VW=(K+1)^{-1},\qquad
 \Delta_KW=-\frac{V}{(K+1)(K+2)}.
\]
The two reciprocal bounds and \(p-s\le(1-\epsilon)p\), \(s\ge\epsilon p\) therefore give
\[
 \operatorname{Var}(W)\le \frac{C_\epsilon}{1+tp}.
\]
Independence of \(S\) and \(W\) now gives the exact identity
\[
 \operatorname{Var}\{(S/u)W\}=\frac q uEW^2+q^2\operatorname{Var}(W).
\]
Summing over cells uses \(\sum_xq_x\le1\) and
\[
 \sum_x\frac{p_x^2}{1+tp_x}\le \frac1t\sum_xp_x=\frac1t.
\]
The arm difference satisfies \(\operatorname{Var}(U-V)\le2\operatorname{Var}(U)+2\operatorname{Var}(V)\), so no arm independence is assumed.  Independent Poisson cells then prove the stated variance bound.

%%% FIELD baseline-statement
There is a finite \(C_\epsilon\), depending only on \(\epsilon\), such that for all \(n\ge1\), \(m\ge0\), and \(d\ge2\),
\[
 R_\epsilon(n,m,d)\le C_\epsilon\min\!\left\{1,\frac1n+\frac{d^2}{(n+m)^2}\right\}.
\]
The witnessing estimator is the clipped and derandomized Poisson-prefix version of the all-heavy inverse-count statistic in \(\mathrm{lem:poisson\mbox{-}inverse\mbox{-}count\mbox{-}risk}\).

%%% FIELD baseline-proof
When \(n\) is below a fixed numerical constant, the zero estimator has risk at most one and proves the assertion after enlarging \(C_\epsilon\).  Otherwise reserve at least \(n/3\) labeled observations for outcome successes and combine the remaining labeled \((X,A)\) projections with all \(m\) auxiliary records; this gives at least \((n+m)/3\) independent marginal records by the three sampling assumptions.  Independently draw Poisson prefix lengths of means one eighth of these two block sizes.  If either length exceeds its block, output zero; otherwise compute \(Z^H\) and clip it to \([-1,1]\).  The two Poisson means satisfy \(u\asymp n\) and \(t\asymp n+m\).  The probability of a prefix overflow is at most \(e^{-c n}+e^{-c(n+m)}\le C/n\), by the exponential-series bound \(P\{\operatorname{Poi}(r/8)>r\}\le(e/8)^{r}\).

On the no-overflow event the selected prefixes have exactly the independent Poisson experiment used in \(\mathrm{lem:poisson\mbox{-}inverse\mbox{-}count\mbox{-}risk}\).  Clipping cannot increase squared error because \(\tau(P)\in[-1,1]\).  Coupling the truncated and untruncated prefix experiments, both with outputs in \([-1,1]\), adds at most four times the overflow probability.  Thus the preceding lemma gives
\[
 \sup_P E[(\widehat\tau-\tau(P))^2]
 \le C_\epsilon\left\{\frac1n+\frac1{n+m}+\frac{d^2}{(n+m)^2}\right\}.
\]
Finally average the randomized output over the finitely many nonoverflow prefix lengths conditional on the data.  Conditional Jensen shows that this deterministic finite weighted sum has no larger risk.  Since \((n+m)^{-1}\le n^{-1}\), the uncapped bound follows, and the zero estimator supplies the cap at one.

%%% FIELD binary-statement
For \(d=2\) there are constants \(0<c_\epsilon\le C_\epsilon<\infty\) such that, uniformly over \(n\ge1\) and \(m\ge0\),
\[
 \frac{c_\epsilon}{n}\le R_\epsilon(n,m,2)\le\frac{C_\epsilon}{n}.
\]

%%% FIELD binary-proof
Supplying the exact table \(P_{XA}\) is more informative than supplying \(m\) draws from it, so every estimator in the latter experiment can be used in the former by simulating and then discarding those draws.  Hence \(R_\epsilon(n,m,2)\ge R_\epsilon^{\mathrm{known}}(n,2)\ge c_\epsilon/n\) by \(\mathrm{thm:known\mbox{-}marginal\mbox{-}boundary}\).  The baseline theorem and \(N=n+m\ge n\) give
\[
 R_\epsilon(n,m,2)\le C_\epsilon\min\{1,n^{-1}+4N^{-2}\}\le5C_\epsilon/n.
\]

%%% FIELD factorial-statement
For every integer \(L\ge2\), the displayed polynomials \(H_L,E_L,G_L\) are well defined, \(G_L\) has degree at most \(L-2\), and on \(0\le z\le1\),
\[
 0\le E_L(z)\le\min\{1,(L^2z)^{-1}\},\qquad zG_L(z)+E_L(z)=1.
\]
There is a universal \(A>1\) such that the factorial lift \(\widehat W_{a,L}\) of \(W_{a,L}\) satisfies, whenever \(tB\ge L\),
\[
 E\widehat W_{a,L}^{,2}\le A^L\{1+p/B\}^{2L},
\]
and the right side can be replaced by \(A^L\) when \(p\le B\).

%%% FIELD factorial-proof
The identity \(T_L(1)=1\) shows that \(H_L\) is divisible by \(z\).  Since \(T_L'(1)=L^2\), its quotient has value one at zero, so \(1-E_L\) is again divisible by \(z\).  This proves the degree and divisibility assertions.  Writing \(1-2z=\cos\theta\) for \(z\in[0,1]\) gives
\[
 E_L(z)=\frac{1-\cos(L\theta)}{L^2(1-\cos\theta)}
 =\left\{\frac{\sin(L\theta/2)}{L\sin(\theta/2)}\right\}^2.
\]
Thus \(0\le E_L\le1\); the numerator bound \(1-\cos(L\theta)\le2\) also gives \(E_L(z)\le(L^2z)^{-1}\).  The identity involving \(G_L\) is its definition.

The recurrence \(T_{j+1}(w)=2wT_j(w)-T_{j-1}(w)\), applied after substituting \(w=1-2z\), shows inductively that the sum of the absolute coefficients of \(T_L(1-2z)\) is at most \(C_0^L\) for a universal \(C_0\).  The two exact divisions only shift coefficients and multiply their sum by a polynomial factor in \(L\), which is absorbed into \(C_1^L\).  Hence the dimensionless polynomial \((z_0+z_1)G_L(z_a)\) has coefficient \(\ell^1\)-norm at most \(C_1^L\).

For \(K\sim\operatorname{Poi}(\lambda)\), direct counting of common ordered selections gives
\[
 E[(K)_r(K)_v]=\sum_{h=0}^{\min(r,v)}\binom rh\binom vh h!\lambda^{r+v-h}.
\]
After division by \((tB)^{r+v}\), each overlap term is bounded by
\[
 \frac{(rv)^h}{h!(tB)^h}\,(1+\lambda/(tB))^{r+v-h}.
\]
Because \(r,v\le L\) and \(tB\ge L\), summing over \(h\) is at most \(\exp(L)\{1+\lambda/(tB)\}^{r+v}\).  Apply this independently to the two arm counts, then use Cauchy--Schwarz for the finite coefficient sum.  Since \((\lambda_0+\lambda_1)/(tB)=p/B\), the result is \(A^L(1+p/B)^{2L}\) after increasing a universal \(A\).  When \(p\le B\), the last factor is at most \(2^{2L}\), which is absorbed into \(A^L\).

%%% FIELD hybrid-statement
There are constants \(H_\epsilon>1\), \(c_\epsilon^\circ>0\), \(n_\epsilon^\circ\), and \(C_\epsilon<\infty\) with the following property.  Let \(n\ge n_\epsilon^\circ\), \(N\ge n\), \(u\asymp n\), and \(t_p\asymp t\asymp N\), with the comparison constants fixed numerically.  Take \(L=\max\{2,\lfloor c_\epsilon^\circ\log(en)\rfloor\}\), \(B=H_\epsilon L/\min(t_p,t)\), an independent pilot count \(J_x\sim\operatorname{Poi}(t_pp_x)\), and the hybrid cell statistic in \(\mathrm{def:mixed\mbox{-}estimator\mbox{-}handle}\).  Then its clipped sum \(Z^{\mathrm{mix}}\) obeys, uniformly over \(P\in\mathcal M_{d,\epsilon}\),
\[
 E_P[(Z^{\mathrm{mix}}-\tau(P))^2]
 \le C_\epsilon\min\!\left\{1,\frac1n+\frac{d^2}{N^2\log^2(en)}\right\}.
\]
The constant is independent of \(d\), \(N/n\), and \(m=N-n\).

%%% FIELD hybrid-proof
Write \(\ell=\log(en)\) and \(b=\min(t_p,t)B=H_\epsilon L\).  All comparison constants below depend only on the fixed ratios \(u/n,t_p/N,t/N\) and on \(\epsilon\).  For cell \(x\), let \(I_x=\mathbf1\{J_x\le\lfloor t_pB/4\rfloor\}\), let \(P_x\) be its polynomial statistic, and let \(H_x\) be its inverse-count statistic.  The pilot, outcome, and factorial blocks are independent, and different cells are independent in the Poisson experiment.

First consider deterministic means.  Put \(z_a=s_{ax}/B\).  The factorial identity and \(zG_L(z)=1-E_L(z)\) give
\[
 E P_x=p_x(\mu_{1x}-\mu_{0x})
 -p_x\{\mu_{1x}E_L(z_1)-\mu_{0x}E_L(z_0)\}.
\]
Consequently, if \(p_x\le B\), overlap and the bound on \(E_L\) imply
\[
 |EP_x-p_x(\mu_{1x}-\mu_{0x})|\le \frac{2B}{\epsilon L^2}. \tag{1}
\]
The inverse-count calculation gives, for every cell,
\[
 |EH_x-p_x(\mu_{1x}-\mu_{0x})|\le2p_xe^{-\epsilon t p_x}. \tag{2}
\]

We next record the tail bounds that repair the false-light mean.  If \(p>B\), the Poisson Chernoff series yields
\[
 \pi_x:=P(I_x=1)\le \exp(-c_0t_pp_x).
\]
The coefficient recurrence in \(\mathrm{lem:chebyshev\mbox{-}factorial\mbox{-}certificate}\) also gives the deterministic continuation bound
\[
 |EP_x|\le p_xA^L(1+p_x/B)^L
\]
after increasing \(A\).  Choose \(H_\epsilon\) so large that, for every \(y\ge1\),
\[
 c_0H_\epsilon y-2\log(1+y)-2\log A\ge c_1H_\epsilon y. \tag{3}
\]
Then choose \(c_\epsilon^\circ>0\) small enough for the moment bounds below, and enlarge \(H_\epsilon\), if necessary, so that \(c_1H_\epsilon L\ge12\ell\).  Equations (3) and \(\sum_xp_x=1\) give directly, without Cauchy--Schwarz,
\[
 \sum_{p_x>B}\pi_x|EP_x|
 \le \sum_{p_x>B}p_xe^{-c_1H_\epsilon Lp_x/B}\le n^{-12}. \tag{4}
\]
Likewise, the second-moment part of the factorial certificate gives
\[
 \sum_{p_x>B}\pi_xEP_x^2\le n^{-10}. \tag{5}
\]
Here a factor \(q_{ax}/u\) or \(q_{ax}^2\) precedes the factorial second moment; both sums are bounded using \(\sum_xq_{ax}\le1\) and \(q_{ax}^2\le p_xq_{ax}\).  Thus neither (4) nor (5) contains a cell count or the ratio \(N/n\).

For the opposite error, put \(k_0=\lfloor t_pB/4\rfloor\).  When \(p\le B\), the exponential-series calculation
\[
 P\{\operatorname{Poi}(z)>k_0\}\le
 \exp\{-k_0\log(k_0/z)+k_0-z\},\qquad z<k_0,
\]
together with the trivial bound one for \(z\ge k_0\), shows that
\[
 P(I_x=0)e^{-\epsilon t p_x}\le e^{-c_2H_\epsilon L}. \tag{6}
\]
Indeed, for \(z=t_pp_x<k_0\), the exponent after multiplying by \(e^{-\epsilon tp_x}\) is maximized at \(z=k_0/(1+\epsilon t/t_p)\) and equals \(-k_0\log(1+\epsilon t/t_p)\); for \(z\ge k_0\), the exponential factor itself proves (6).  Increase \(H_\epsilon\) so that (6) is at most \(n^{-12}\).  Combining (1), (2), (4), and (6), and summing with \(\sum_xp_x=1\), gives
\[
 |E\sum_x\{I_xP_x+(1-I_x)H_x\}-\tau(P)|
 \le C_\epsilon\frac{dB}{L^2}+Cn^{-10}. \tag{7}
\]
This is the direct expectation calculation; no square root of \(p_x/u\) is introduced.

At the requested early-kill value \(\epsilon=1/4\), (4) is the bound on the false-light side \(p_x>B\), while (6) is the bound on the false-heavy side \(p_x\le B\).  Both are uniform at the threshold because \(t_pB\asymp L\), and both are summed with \(p_x\), not with the number of cells.  Thus the test produces neither \(d e^{-cL}\) nor a factor involving \(N/n\).

Conditional on the pilot,
\[
 \operatorname{Var}(I_xP_x+(1-I_x)H_x)
 =\pi_x\operatorname{Var}(P_x)+(1-\pi_x)\operatorname{Var}(H_x)
 +\pi_x(1-\pi_x)(EP_x-EH_x)^2. \tag{8}
\]
For \(p_x\le B\), independence of the success count and factorial lift gives, arm by arm,
\[
 \operatorname{Var}\{(S_{ax}/u)\widehat W_{a,L,x}\}
 =\frac{q_{ax}}uE\widehat W_{a,L,x}^2
 +q_{ax}^2\operatorname{Var}(\widehat W_{a,L,x}).
\]
The factorial certificate, \(q_{ax}\le p_x\), and \(p_x\le B\) yield
\[
 \sum_{p_x\le B}\operatorname{Var}(P_x)
 \le C A^L\left\{\frac{\min(1,dB)}u+dB^2\right\}. \tag{9}
\]
The arm covariance is covered by \(\operatorname{Var}(U-V)\le2\operatorname{Var}(U)+2\operatorname{Var}(V)\).  For \(p_x>B\), (5) bounds the false-light contribution.  The all-heavy lemma bounds the sum of the inverse-count variances by \(C_\epsilon(u^{-1}+t^{-1})\).

For the random branch means in (8), (1)--(2) imply, when \(p_x\le B\),
\[
 (EP_x-EH_x)^2\le C_\epsilon\{B^2/L^4+p_x^2e^{-2\epsilon tp_x}\}.
\]
Their sum is at most
\[
 C_\epsilon\{dB^2/L^4+t^{-1}\}, \tag{10}
\]
because \(p^2e^{-2\epsilon tp}\le C_\epsilon p/t\).  For \(p_x>B\), (3)--(5) and \(|EH_x|\le C_\epsilon p_x\) give \(n^{-10}\).  Equations (8)--(10), independence across cells, and (7) give
\[
 E[(Z^{\mathrm{mix}}-\tau(P))^2]
 \le C_\epsilon\left[
 \frac1n+\frac1N+\frac{d^2B^2}{L^4}
 +\frac{A^L\min(1,dB)}n+A^LdB^2+\frac{dB^2}{L^4}+n^{-9}
 \right]. \tag{11}
\]

Set \(\delta=d/(NL)\).  Since \(B\asymp L/N\), choose \(c_\epsilon^\circ\) small enough that
\[
 L^4A^{2L}\le n,
 \qquad L^3A^L\le\sqrt n. \tag{12}
\]
If \(\delta L^2\le1\), the product inequality and (12) give \(A^LdB/n\le C(\delta^2+n^{-1})\); if \(\delta L^2>1\), then \(A^L/n\le L^{-4}\le\delta^2\).  Thus the term with \(\min(1,dB)\) has the same bound.  Moreover,
\[
 A^LdB^2\le C\frac{\delta L^3A^L}{N}\le C(\delta^2+n^{-1}),
 \qquad
 \frac{dB^2}{L^4}\le C\frac{\delta}{NL}\le C(\delta^2+n^{-1}).
\]
Substitution into (11) proves the uncapped rate.  Projection to \([-1,1]\) cannot increase loss, and its trivial risk bound is four; enlarging \(C_\epsilon\) supplies the cap and covers the finitely many \(n<n_\epsilon^\circ\).

%%% FIELD upper-transfer-statement
The independent-Poisson hybrid statistic of \(\mathrm{lem:poisson\mbox{-}hybrid\mbox{-}risk}\) has a total deterministic implementation from exactly \(n\) labeled and \(m\) auxiliary observations.  Its risk is at most twice the Poisson risk plus \(C/n\), uniformly over the observation experiment.

%%% FIELD upper-transfer-proof
For \(n\) above a fixed constant, partition the labeled indices into three deterministic blocks of sizes between \(n/4\) and \(n/2\).  The first is the outcome block.  Concatenate the \((X,A)\) projections from each of the other two with disjoint halves of the auxiliary sample; these are independent identically distributed \(P_{XA}\) observations by the three sampling assumptions, and each resulting marginal block has size at least \(N/4-1\).  Independently draw three Poisson prefix lengths with means one eighth of their block sizes.  Their means satisfy \(u\asymp n\) and \(t_p\asymp t\asymp N\).  On overflow output zero; otherwise use the prefixes in the hybrid construction and clip.

Extend each finite block by an independent infinite sequence from its own law.  Under this coupling the valid-prefix output equals the untruncated Poisson statistic.  The union overflow probability is bounded by \(Ce^{-cn}\), since \(N\ge n\).  Both coupled outputs lie in \([-1,1]\), so
\[
 E[(Z_{\rm trunc}-\tau)^2]\le2E[(Z_{\rm Pois}-\tau)^2]+8Ce^{-cn}.
\]
The auxiliary prefix variables take only finitely many nonoverflow values in the displayed output, the overflow value being zero.  Define the reported estimator as \(E[Z_{\rm trunc}\mid\mathcal L_n,\mathcal U_m]\), an explicit finite sum with Poisson weights.  Conditional Jensen makes its risk no larger.  Finally \(e^{-cn}\le C/n\).  For the finitely many smaller \(n\), the zero estimator is the stipulated total branch.

%%% FIELD upper-statement
With
\[
 r_\epsilon(n,m,d)=\min\!\left\{1,\frac1n+\frac{d^2}{(n+m)^2\log^2(en)}\right\},
\]
the total computable estimator \(\widehat\tau^{\mathrm{mix}}_{n,m,d}\) in \(\mathrm{def:mixed\mbox{-}estimator\mbox{-}handle}\) satisfies, for a constant depending only on \(\epsilon\),
\[
 \sup_{P\in\mathcal M_{d,\epsilon}}
 E_{P^{\otimes n}\otimes P_{XA}^{\otimes m}}
 [\{\widehat\tau^{\mathrm{mix}}_{n,m,d}-\tau(P)\}^2]
 \le C_\epsilon r_\epsilon(n,m,d)
\]
simultaneously for every \(n\ge1\), \(m\ge0\), and \(d\ge2\).

%%% FIELD upper-proof
The estimator construction partitions the observed records into independent blocks permitted by \(\mathrm{ass:labeled\mbox{-}iid}\), \(\mathrm{ass:auxiliary\mbox{-}iid}\), and \(\mathrm{ass:sample\mbox{-}independence}\).  The valid prefixes therefore have the exact Poisson laws in \(\mathrm{lem:poisson\mbox{-}hybrid\mbox{-}risk}\), which proves the claimed rate there without a dimension, \(m/n\), or \(\log(n+m)\) remainder.  The deterministic fixed-sample implementation and its \(O(n^{-1})\) coupling cost are proved in \(\mathrm{lem:fixed\mbox{-}sample\mbox{-}hybrid\mbox{-}transfer}\).  Since \(n^{-1}\le r_\epsilon(n,m,d)\), that cost is absorbed.  The small-sample zero branch and clipping give the capped assertion.  All thresholds, degrees, coefficients, and prefix weights are functions only of \(n,m,d,\epsilon\), so the map is total, computable, and contains no unknown-law optimization.

%%% FIELD duality-statement
Fix \(0<\epsilon<1/2\) and put \(\kappa=(1-2\epsilon)/\epsilon\).  There are positive constants \(\gamma_\epsilon,c_\epsilon\) such that, for every integer \(L\ge2\), every \(B>0\), and \(a=\gamma_\epsilon B/L^2\), the interval \(I=[a/\kappa,B]\) is nonempty and there is a finite signed measure \(\sigma\) on \(I\) satisfying
\[
 |\sigma|(I)=1,\qquad \int_Ip^j\,d\sigma(p)=0\quad(0\le j\le L),
 \qquad \int_I\frac p{p+a}\,d\sigma(p)\ge c_\epsilon.
\]

%%% FIELD duality-proof
Let \(f(p)=p/(p+a)\).  The points \(p_1=a/\kappa\) and \(p_2=2a/\kappa\) belong to \(I\) once \(\gamma_\epsilon\le\kappa/2\), and
\[
 f(p_2)-f(p_1)=\frac{\kappa}{(1+\kappa)(2+\kappa)}=:D_\epsilon>0. \tag{13}
\]
Suppose a polynomial \(Q\) of degree at most \(L\) satisfies \(\|f-Q\|_{\infty,I}<D_\epsilon/4\).  Then the mean value theorem and (13) give a point \(\xi\in(p_1,p_2)\) for which
\[
 |Q'(\xi)|\ge \frac{D_\epsilon\kappa}{2a}. \tag{14}
\]
Also \(\|Q\|_{\infty,I}\le1+D_\epsilon/4\).  The polynomial derivative inequality
\[
 \|Q'\|_{\infty,[r,s]}\le \frac{2L^2}{s-r}\|Q\|_{\infty,[r,s]} \tag{15}
\]
follows as follows.  After affine reduction to \([-1,1]\), interpolate \(Q\) at the \(L+1\) extrema of \(T_L\); differentiating the Lagrange formula and pairing symmetric nodes bounds the absolute sum of the derivative cardinal functions by \(L^2\).  The endpoint value is exactly \(L^2\), and applying the same formula after moving the evaluation point between its two adjacent extrema gives no larger sum.  Rescaling proves (15).  Applying (15) to \(I\), then choosing
\[
 0<\gamma_\epsilon<\frac{D_\epsilon\kappa}{8(1+D_\epsilon/4)},
\]
contradicts (14).  Hence the best uniform approximation error of \(f\) by degree-\(L\) polynomials is at least \(D_\epsilon/4\).

We now give the finite dual measure rather than invoke it.  A best approximating polynomial \(Q_*\) exists: a minimizing sequence is bounded in coefficient norm because all norms on the finite-dimensional polynomial space are equivalent, and hence has a convergent subsequence.  If the error \(f-Q_*\) did not attain alternating values \(\eta(-1)^iE\), \(E=\|f-Q_*\|_\infty\), at \(L+2\) ordered points, its extremal set would have at most \(L+1\) alternating sign blocks.  Put one root between consecutive blocks and orient the resulting degree-at-most-\(L\) polynomial \(R\) to have the error's sign on every block.  Compactness gives a neighborhood of the extremal set on which the signs agree and a strict error margin off that neighborhood; therefore \(\|f-(Q_*+tR)\|_\infty<E\) for all sufficiently small \(t>0\), a contradiction.  Thus such points \(z_0<\cdots<z_{L+1}\) exist.

Let
\[
 w_i=\left|\prod_{j\ne i}(z_i-z_j)\right|^{-1},\qquad
 \sigma_0=\eta\sum_{i=0}^{L+1}(-1)^iw_i\delta_{z_i}.
\]
The coefficient of \(z^{L+1}\) in the Lagrange interpolation formula shows that, for every polynomial \(R\) of degree at most \(L\),
\[
 \sum_i\frac{R(z_i)}{\prod_{j\ne i}(z_i-z_j)}=0.
\]
Since the signs of the denominators alternate, \(\sigma_0\) annihilates all such \(R\).  Normalize it by \(\sum_iw_i\).  Its total variation is one, and
\[
 \int f\,d\sigma=\int(f-Q_*)\,d\sigma=E\ge D_\epsilon/4.
\]
Taking \(c_\epsilon=D_\epsilon/4\) proves the claim.

%%% FIELD prior-statement
There are constants \(b_\epsilon,\gamma_\epsilon,c_\epsilon,C_\epsilon>0\) such that the finite-atom construction in \(\mathrm{def:common\mbox{-}marginal\mbox{-}prior\mbox{-}handle}\), with \(L\ge2\), \(B=b_\epsilon L/M\), \(a=\gamma_\epsilon B/L^2\), and \(k\) rare cells, has the following properties whenever \(ka\le c_\epsilon\): it lies in \(\mathcal M_{d,\epsilon}\), its two priors have exactly the same distribution of the full random table \(P_{XA}\), its raw target centers are separated by at least \(c_\epsilon ka\), its raw mass and target variances are at most \(C_\epsilon kBa\), and the total variation distance of the two mixture laws in a Poisson experiment of labeled intensity \(u\) and auxiliary intensity \(v\), \(M=u+v\), is at most
\[
 C_\epsilon uka\rho_\epsilon^L
\]
for a fixed \(0<\rho_\epsilon<1\).  The likelihood discrepancy is gated by the labeled intensity \(u\), although \(M\) includes all auxiliary observations.

%%% FIELD prior-proof
Take \(\sigma\) from \(\mathrm{lem:finite\mbox{-}alternation\mbox{-}duality}\), write \(h=d\sigma/d|\sigma|\), and define \(\nu\) as in the construction.  Its atom at zero is nonnegative because \(a/(p+a)\le1\).  Moreover,
\[
 E_\nu p=a\int_I\frac p{p+a}\,d|\sigma|(p)\le a,
 \qquad E_\nu p^2\le B E_\nu p\le Ba. \tag{16}
\]
For \(p>0\), set \(s=\epsilon(p+a)\).  Since \(p\ge a/\kappa\),
\[
 \epsilon\le e(p)=\epsilon(1+a/p)\le\epsilon(1+\kappa)=1-\epsilon. \tag{17}
\]
The reservoir mass \(p_*=1-kE_\nu p\) is positive when \(ka\le c_\epsilon<1/2\).  The raw total \(S=p_*+\sum_ip_i\) is therefore positive almost surely, and division by \(S\) produces nonnegative masses summing to one without changing any propensity.  Equations (17) and the identical legal reservoir show membership in \(\mathcal M_{d,\epsilon}\).  The latent vector \((p_i,e_i)_{i\le k}\), the reservoir, and \(S\) have the same joint distribution under both hypotheses; hence the normalized full table \(P_{XA}\) has exactly the same distribution, not merely matching moments.

Let \(T_j=\sum_ip_i\mu_{1i,j}\) be the raw treated target under prior \(j\).  Because \(\mu_{1i,1}-\mu_{1i,0}=h(p_i)\),
\[
 E(T_1-T_0)=k\int p h(p)\,d\nu(p)
 =ka\int\frac p{p+a}\,d\sigma(p)\ge c_\epsilon ka. \tag{18}
\]
Independence, (16), and \(0\le p\mu\le p\) give
\[
 \operatorname{Var}(S)+\max_j\operatorname{Var}(T_j)\le2kBa. \tag{19}
\]

It remains to verify the joint likelihood.  In one rare cell of the unnormalized Poisson experiment, aggregate the labeled treated count as \(R\), auxiliary treated count as \(U\), and all control counts as \(C\).  Conditional splitting of controls into labeled and auxiliary parts is parameter-free.  If \(R=0\), the two mixture laws agree.  For \(r\ge1,j,k\ge0\), the difference for either nonzero labeled outcome mark is
\[
 \int h(p)e^{-Mp}\frac{(us)^r(vs)^j\{M(p-s)\}^k}{r!j!k!}\,d\nu(p).
\]
The two possible labeled marks cancel the factor one half in total variation.  Factor one \(us\) and use
\[
 h(p)s\,d\nu(p)=\epsilon a\,d\sigma(p). \tag{20}
\]
After expanding \(e^{-Mp}\), every term whose remaining total polynomial degree is at most \(L\) vanishes by the moment equations for \(\sigma\).  Absolute convergence is justified by \(p\le B\), so Tonelli applies after taking the absolute majorant.  Summing only the nonvanishing degrees and using \(s\le p\), \(p-s\le p\), and \(u+v=M\), gives
\[
 \operatorname{TV}_{\rm cell}
 \le u\epsilon a\sum_{r>L}\frac{(2MB)^r}{r!}. \tag{21}
\]
Choose \(b_\epsilon\) so small that \(2eMB/L=2eb_\epsilon<1/2\).  The ratio test then bounds (21) by \(C_\epsilon ua\rho_\epsilon^L\) for a fixed \(\rho_\epsilon<1\).  The rare cells are independent in the raw Poisson experiment, while the reservoir has a common law.  The product inequality \(\operatorname{TV}(\otimes_iQ_i,\otimes_iR_i)\le\sum_i\operatorname{TV}(Q_i,R_i)\), obtained by replacing one factor at a time, proves the stated \(C_\epsilon uka\rho_\epsilon^L\) bound.  This establishes both label gating and the absence of an auxiliary-only discrepancy.

%%% FIELD lower-transfer-statement
Let the priors of \(\mathrm{lem:common\mbox{-}marginal\mbox{-}poisson\mbox{-}prior}\) have raw center separation \(\Delta\), mixture total variation at most \(1/8\), and satisfy \(kBa\le c\Delta^2\) for a sufficiently small numerical \(c\).  If their Poisson intensities obey \(u=Cn\) and \(v=C(n+m)\) for a sufficiently large fixed \(C\), then
\[
 R_\epsilon(n,m,d)\ge c'\Delta^2.
\]
The transfer preserves the identical random \(P_{XA}\) law and holds for the fixed, unequal sample sizes.

%%% FIELD lower-transfer-proof
Let \(S\) and \(T_j\) be the raw total mass and raw target.  The normalized target is \(\tau_j=T_j/S\), and \(0\le T_j\le S\), so
\[
 |\tau_j-E T_j|\le |T_j-E T_j|+|S-1|. \tag{22}
\]
By (19), Chebyshev's inequality and \(kBa\le c\Delta^2\) imply
\[
 \Pi_j\{|\tau_j-E T_j|>\Delta/4\}\le32kBa/\Delta^2\le1/16 \tag{23}
\]
after fixing \(c\).  Thresholding any estimator halfway between the two centers gives a test.  If its Bayes mean squared error under the two priors were below \(\Delta^2/256\), Markov's inequality and (23) would make the sum of its two testing errors smaller than \(1/2\).  But every test has error sum at least \(1-\operatorname{TV}\ge7/8\).  Hence the Bayes risk in the raw Poisson experiment is at least \(c_0\Delta^2\).

We next transfer an arbitrary fixed-sample estimator \(T\), clipped to \([-1,1]\), to that raw experiment.  Given the raw cell masses, the labeled Poisson process has total count \(K_L\sim\operatorname{Poi}(uS)\); conditional on \(K_L\), its ordered records are independent from the normalized law.  The same statement holds independently for the auxiliary process, with \(K_U\sim\operatorname{Poi}(vS)\).  On \(K_L\ge n,K_U\ge m\), apply \(T\) to the first \(n\) and \(m\) records and discard the rest.  This conditional experiment is exactly \(P^{\otimes n}\otimes P_{XA}^{\otimes m}\).

On \(S\ge3/4\), choosing the numerical \(C\) large enough gives
\[
 P(K_L<n)+P(K_U<m)\le e^{-c_1n}+e^{-c_1(n+m)}.
\]
On the complementary event, Chebyshev and (19) give \(P(S<3/4)\le16kBa\).  In the regime where this lemma is used, \(\Delta^2\ge c_2/n\), so the exponential remainder is at most \(C\Delta^2\); choosing the earlier concentration constant and then \(C\) makes the complete failure probability at most \(c_0\Delta^2/8\).  Output zero on failure.  Since squared loss is at most four after clipping, the transferred procedure has Bayes risk at most
\[
 \sup_{P\in\mathcal M_{d,\epsilon}}E_P[(T-\tau(P))^2]+c_0\Delta^2/2.
\]
Comparison with the raw Bayes lower bound yields the claim after taking the infimum over \(T\).  The normalized priors are the same priors throughout this transfer, so their full random \(P_{XA}\) law remains identical.

%%% FIELD lower-statement
For the explicit rate
\[
 r_\epsilon(n,m,d)=\min\!\left\{1,\frac1n+\frac{d^2}{(n+m)^2\log^2(en)}\right\},
\]
the regime-dependent priors \(\Pi_0,\Pi_1\) in \(\mathrm{def:common\mbox{-}marginal\mbox{-}prior\mbox{-}handle}\) induce exactly the same distribution of the entire random table \(P_{XA}\) and imply, for a constant depending only on \(\epsilon\),
\[
 R_\epsilon(n,m,d)\ge c_\epsilon r_\epsilon(n,m,d)
\]
for every \(n\ge1\), \(m\ge0\), and \(d\ge2\), including the constant-risk regime.

%%% FIELD lower-proof
Put \(N=n+m\) and \(\ell=\log(en)\).  For the parametric branch, concentrate \(X\) on one cell, take \(e=1/2\), \(\mu_0=0\), and compare \(\mu_1=1/2-h_n\) with \(\mu_1=1/2+h_n\), where \(h_n=c_0/\sqrt n\le1/4\).  The two auxiliary laws coincide.  For one labeled observation, direct summation over its four possible \((A,Y)\) values gives chi-squared divergence at most \(C h_n^2\); independence gives product divergence at most \((1+Ch_n^2)^n-1\le e^{Cc_0^2}-1\).  Choose \(c_0\) so this is at most \(1/4\).  Total variation is then at most one half by Cauchy--Schwarz, while the target separation is \(2h_n\).  Thresholding any estimator at the midpoint and using the same squared-loss testing argument as in \(\mathrm{lem:fixed\mbox{-}sample\mbox{-}common\mbox{-}marginal\mbox{-}transfer}\) yields
\[
 R_\epsilon(n,m,d)\ge c_\epsilon/n. \tag{24}
\]

It remains to obtain the dimension term when it is larger than a sufficiently large multiple of \(n^{-1}\).  The finitely many \(n<n_\epsilon^\circ\) are already covered by (24), after reducing the final constant, because the desired capped rate is at most one.  For larger \(n\), take
\[
 L=\lceil C_0\ell\rceil,\qquad
 M=C_1n+C_1N,\qquad
 B=b_\epsilon L/M,\qquad
 a=\gamma_\epsilon B/L^2,
\]
and
\[
 k=\min\{d-1,\lfloor c_0NL\rfloor\}.
\]
The constants \(b_\epsilon,\gamma_\epsilon\) are those of the common-marginal prior lemma, \(C_0\) is large enough that \(C_\epsilon n\rho_\epsilon^L\le1/8\), \(C_1\) is the fixed transfer intensity, and \(c_0\) is small enough that \(ka<1/4\).  Since \(M\asymp N\), the center separation (18) satisfies
\[
 \Delta\ge c_\epsilon ka\ge c_\epsilon\min\!\left\{1,\frac d{N\ell}\right\}. \tag{25}
\]
The replacement of \(d-1\) by \(d\) costs only a factor two because \(d\ge2\).

Suppose first that the square of the right side of (25) is at least \(C_2/n\).  Then
\[
 k\ge c_\epsilon\sqrt n\,\ell,
\]
whether or not the cap in (25) is active.  As \(\sqrt n/\ell\to\infty\), increasing \(n_\epsilon^\circ\) and then \(C_2\) ensures \(k\ge C_3L^2\).  From \(B/a=L^2/\gamma_\epsilon\),
\[
 \frac{kBa}{\Delta^2}\le C_\epsilon\frac{L^2}{k}\le c,
\]
where \(c\) is the transfer constant.  The prior lemma gives mixture total variation at most \(1/8\), because \(uka\rho_epsilon^L\le Cn
ho_epsilon^L\).  Thus \(\mathrm{lem:fixed\mbox{-}sample\mbox{-}common\mbox{-}marginal\mbox{-}transfer}\) and (25) yield
\[
 R_\epsilon(n,m,d)\ge c_\epsilon\min\!\left\{1,\frac{d^2}{N^2\ell^2}\right\}. \tag{26}
\]
If the dimension term is smaller than \(C_2/n\), (24) already dominates it.  Combining (24) and (26), using \(\max(r,s)\ge(r+s)/2\) before the outer cap, proves the displayed lower rate in all regimes.  Membership, normalization, exact auxiliary invariance, label-gated likelihood cancellation, tensorization, and the fixed unequal-sample transfer are respectively equations (16)--(21) and the transfer lemma; no premise about a minimum cell mass was used.

%%% FIELD known-limit-statement
For every fixed \(n\ge1\) and \(d\ge2\),
\[
 \lim_{m\to\infty}R_\epsilon(n,m,d)=R_\epsilon^{\mathrm{known}}(n,d).
\]

%%% FIELD known-limit-proof
Write the observed law by its finite probability vector and let \(Q=P_{XA}\).  The model is a compact subset of a finite-dimensional simplex: overlap is the pair of closed inequalities \(s_{1x}\ge\epsilon p_x\) and \(s_{0x}\ge\epsilon p_x\), and the binary outcome masses satisfy closed simplex inequalities.  The map
\[
 P\longmapsto\tau(P)=\sum_x\left(\frac{p_xq_{1x}}{s_{1x}}-\frac{p_xq_{0x}}{s_{0x}}\right),
\]
with each null-cell summand set to zero, is continuous because overlap bounds its absolute value by \(2p_x\) near a null cell.  The labeled sample space is finite for fixed \(n,d\), so risks of bounded decision rules are continuous in the law.

Exact knowledge of \(Q\) can simulate \(m\) auxiliary observations.  Conditional averaging over that simulation cannot increase squared loss, so
\[
 R_\epsilon(n,m,d)\ge R_\epsilon^{\mathrm{known}}(n,d). \tag{27}
\]
For the reverse limit, fix \(\eta>0\).  Partition the compact set of possible \(Q\)'s into finitely many sets of diameter at most \(\delta\), and enlarge each set to its closed \(2\delta\)-neighborhood.  For every enlarged set choose a bounded labeled-sample rule whose maximum risk over laws with marginal in that neighborhood is within \(\eta\) of that local minimax value.  Such rules exist by the definition of the finite-dimensional infimum; projection to \([-1,1]\) is harmless.  Compactness and continuity imply, by contradiction along a convergent sequence as \(\delta\downarrow0\), that the largest of these local values decreases to
\[
 \sup_Q\inf_T\sup_{P:P_{XA}=Q}E_P[(T-\tau(P))^2]
 =R_\epsilon^{\mathrm{known}}(n,d). \tag{28}
\]
The equality holds because the exact table is observed and therefore permits a separate rule on each fiber.

Let \(\widehat Q_m\) be the empirical auxiliary table and use the rule attached to the partition cell containing \(\widehat Q_m\).  If \(\|\widehat Q_m-Q\|_1<\delta\), the true law belongs to that cell's enlarged neighborhood.  Since there are \(2d\) categories, Hoeffding's inequality followed by a union bound gives, uniformly in \(Q\),
\[
 P_Q(\|\widehat Q_m-Q\|_1\ge\delta)
 \le4d\exp\{-m\delta^2/(2d^2)\}\longrightarrow0. \tag{29}
\]
The loss is at most four.  Equations (28)--(29) therefore show \(\limsup_mR_\epsilon(n,m,d)\le R_\epsilon^{\mathrm{known}}(n,d)+2\eta\).  Letting \(\eta\downarrow0\) and using (27) proves the limit.

%%% FIELD frontier-statement
Uniformly over every \(n\ge1\), \(m\ge0\), and \(d\ge2\),
\[
 c_\epsilon\min\!\left\{1,\frac1n+\frac{d^2}{(n+m)^2\log^2(en)}\right\}
 \le R_\epsilon(n,m,d)\le
 C_\epsilon\min\!\left\{1,\frac1n+\frac{d^2}{(n+m)^2\log^2(en)}\right\}.
\]
At \(m=0\) this is
\[
 r_\epsilon(n,0,d)=\min\!\left\{1,\frac1n+\frac{d^2}{n^2\log^2(en)}\right\},
\]
the endpoint recorded in \(\texttt{stat\_discrete\_ate\_minimax\_loggap}\).  For every fixed \(n,d\),
\[
 \lim_{m\to\infty}R_\epsilon(n,m,d)=R_\epsilon^{\mathrm{known}}(n,d).
\]
Along arbitrary sequences \(n\to\infty\), with \(N_n=n+m_n\) and \(\ell_n=\log(en)\), the following are necessary and sufficient:
\[
 R_\epsilon(n,m_n,d_n)\to0
 \quad\Longleftrightarrow\quad
 \frac{d_n}{N_n\ell_n}\to0;
\]
\[
 R_\epsilon(n,m_n,d_n)=O(n^{-1})
 \quad\Longleftrightarrow\quad
 d_n=O\!\left(\frac{N_n\ell_n}{\sqrt n}\right);
\]
and
\[
 R_\epsilon(n,m_n,d_n)=o\{R_\epsilon(n,0,d_n)\}
\]
if and only if both
\[
 \frac{d_n}{\sqrt n\,\ell_n}\to\infty
 \quad\text{and}\quad
 \frac{d_n/(N_n\ell_n)}{\min\{1,d_n/(n\ell_n)\}}\to0.
\]

%%% FIELD frontier-proof
The two finite-sample inequalities are exactly \(\mathrm{thm:uniform\mbox{-}mixed\mbox{-}upper}\) and \(\mathrm{thm:common\mbox{-}marginal\mbox{-}converse}\).  Setting \(m=0\) makes \(N=n\) and gives the displayed supervised expression algebraically; the external record \(\texttt{stat\_discrete\_ate\_minimax\_loggap/oeq:sharp\mbox{-}minimax}\) supplies its prior attribution, not a premise of the present proof.  The exact fixed-parameter assertion is \(\mathrm{thm:known\mbox{-}marginal\mbox{-}limit}\), while \(\mathrm{thm:known\mbox{-}marginal\mbox{-}boundary}\) identifies its order as \(n^{-1}\).

It remains only to verify the sequence characterizations.  Since \(n^{-1}\to0\), the numerical rate tends to zero exactly when \(d_n/(N_n\ell_n)\to0\).  It is \(O(n^{-1})\) exactly when \(d_n^2/(N_n^2\ell_n^2)=O(n^{-1})\), which is the second displayed condition.  The two-sided constants transfer both equivalences from the numerical rate to minimax risk.

For strict improvement, put \(a_n=d_n/(n\ell_n)\) and \(b_n=d_n/(N_n\ell_n)\).  The supervised and semi-supervised numerical rates are respectively
\[
 r_{0n}=\min\{1,n^{-1}+a_n^2\},\qquad
 r_{mn}=\min\{1,n^{-1}+b_n^2\}.
\]
The ratio can tend to zero only if \(n^{-1}=o(r_{0n})\).  Directly from the first formula, this is equivalent to \(a_n\sqrt n=d_n/(\sqrt n\ell_n)\to\infty\).  Under that condition,
\[
 r_{0n}\asymp\min\{1,a_n^2\},
\]
and hence \(r_{mn}=o(r_{0n})\) is equivalent to
\[
 b_n=o\{\min(1,a_n)\}.
\]
Writing this last relation without shorthand gives the second stated limit.  Conversely, the two limits make both \(n^{-1}/r_{0n}\) and \(b_n^2/r_{0n}\) tend to zero, so \(r_{mn}/r_{0n}\to0\).  Applying the uniform two-sided minimax constants at \(m=m_n\) and at \(m=0\) completes necessity and sufficiency.

%%% FIELD prose-tldr
With \(n\) chart-reviewed outcomes and \(m\) additional treatment--covariate records, the unrestricted discrete ATE has minimax mean squared error
\[
 R_\epsilon(n,m,d)\asymp_\epsilon
 \min\!\left\{1,\frac1n+\frac{d^2}{(n+m)^2\log^2(en)}\right\}.
\]
The attaining estimator uses one labeled outcome factor in each light-cell polynomial monomial and learns every remaining treatment--mass factor from the pooled records.  A normalized fuzzy prior keeps the full random \(P_{XA}\) table identically distributed while changing only conditional outcome marks, so auxiliary-only observations cannot erase the \(n^{-1}\) floor.  The result recovers the supervised endpoint, converges for fixed \((n,d)\) to the exact known-\(P_{XA}\) minimax risk, and gives necessary and sufficient annotation-growth thresholds.

%%% FIELD prose-gap
Existing discrete-ATE theory establishes the labeled-only logarithmic rare-cell obstruction, while regular semi-supervised causal theory works under nuisance consistency, working models, or fixed-complexity conditions.  Neither gives a uniform finite-sample theorem when outcome-bearing moments have only \(n\) observations but treatment--mass moments have \(n+m\), null cells are allowed, and the outcome regressions are unrestricted.

%%% FIELD prose-niche
The missing object is the unequal-information annotation frontier.  Its upper proof must prevent false-light polynomial continuation from creating a dimension or \(m/n\) remainder, and its converse must preserve the entire random treatment--covariate law rather than import a supervised prior that auxiliary observations can detect.

%%% FIELD prose-fill
The note supplies the matched rate \(\min\{1,n^{-1}+[d/\{(n+m)\log(en)\}]^2\}\), a total derandomized hybrid estimator, and a normalized common-marginal fuzzy converse.  It also proves the independent inverse-count baseline \(R_\epsilon(n,m,d)\le C_\epsilon\min\{1,n^{-1}+d^2/(n+m)^2\}\), the uniform \(d=2\) parametric rate, the exact fixed-\((n,d)\) known-marginal limit, and elementary necessary-and-sufficient conditions for consistency, root-\(n\) MSE, and strict order improvement.

%%% FIELD prose-related
Zeng, Balakrishnan, Han, and Kennedy, \cite{ZengBalakrishnanHanKennedy2024}, analyze the supervised high-dimensional discrete-ATE obstruction; their observable class agrees with the present class modulo irrelevant null-cell propensities, while their stated Theorem 2 regime does not replace the internally catalogued all-regime supervised endpoint.  Cheng, Ananthakrishnan, and Cai, \cite{ChengAnanthakrishnanCai2021}, Sections 2.1 and 4, motivate random chart-review labeling and give robust locally efficient semi-supervised procedures using additional structure and outcome surrogates.  Chakrabortty and Dai, \cite{ChakraborttyDai2024}, and Kato, \cite{Kato2025}, give adjacent regular asymptotic theories under nuisance-rate conditions.  Jiao, Venkat, Han, and Weissman, \cite{JiaoVenkatHanWeissman2015}, and Wu and Yang, \cite{WuYang2016}, provide the classical polynomial-estimation context.  The present proof's new steps are the unequal-sample factorial lift, the mass-weighted false-branch analysis, and a prior whose full random \(P_{XA}\) law agrees exactly.  The \(m=0\) expression is attributed to \(\texttt{stat\_discrete\_ate\_minimax\_loggap}\); it is an endpoint check rather than the new contribution.

%%% FIELD prose-interpretation
Consistency requires \(d=o\{(n+m)\log(en)\}\), while parametric mean squared error requires and is guaranteed by \(d=O\{(n+m)\log(en)/\sqrt n\}\).  Thus inexpensive treatment--covariate records can remove the high-cardinality approximation term, but they cannot improve the irreducible \(n^{-1}\) cost of outcome annotation.  Strict order improvement over labeled-only estimation additionally requires that the supervised problem be genuinely nonparametric and that the auxiliary sample shrink the approximation term relative to its supervised value.

%%% FIELD prose-limitation
The theorem is minimax over an unrestricted finite alphabet and hides constants depending on the fixed overlap level \(\epsilon\); the construction is intended as a sharp benchmark rather than a computationally optimized implementation.  The fixed-\((n,d)\) limit is an equality of minimax risks, whereas the uniform finite-sample theorem determines those risks only up to \(\epsilon\)-dependent constants.  No claim is made for selected enrollment, covariate shift, post-treatment surrogates, continuous covariates, or violations of overlap.

%%% FIELD prose-scope
The result applies to independent random outcome labeling from one shared population, binary treatment and outcome, known finite alphabet size, occupied-cell overlap, and otherwise unrestricted cell masses, propensities, and outcome regressions.  It gives a baseline-only chart-review allocation benchmark for the Cheng--Ananthakrishnan--Cai setting.  It neither claims their surrogate-based efficiency gain nor asserts that their empirical data obey this unrestricted discrete model.

%%% FIELD note-upper-justification
This theorem isolates and closes the unequal-information upper step: every outcome-bearing monomial uses a labeled factor, while all higher factorial powers use treatment--covariate records from the pooled budget.

%%% FIELD note-upper-gap
The direct deterministic-mean calculation on false-light cells, followed by a separate second-moment argument, removes the square-root leakage that would otherwise introduce \(m/n\); mass-weighted summation removes a hidden \(d e^{-cL}\) remainder.

%%% FIELD note-upper-consumer
The estimator is the model-free discrete upper benchmark for allocating additional chart-review outcomes versus treatment--covariate-only records.

%%% FIELD note-lower-justification
The converse makes auxiliary invariance exact at the level of the full random \(P_{XA}\) table and gates every mixture-likelihood discrepancy by at least one labeled outcome mark.

%%% FIELD note-lower-gap
The normalization is transferred through an unnormalized Poisson experiment whose ordered records have the normalized law conditional on total counts; this controls the reservoir cost without a factor depending on \(m\).

%%% FIELD note-lower-consumer
The lower bound prevents planning rules from counting treatment--covariate-only records as if they also supplied outcome-bearing moments.

%%% FIELD note-frontier-justification
The matched theorem quantifies exactly when auxiliary records change the minimax order while retaining the labeled-only and known-marginal boundaries.

%%% FIELD note-frontier-gap
The rate is not obtained by substituting \(n+m\) for \(n\): only the approximation term uses \(n+m\), whereas the outcome-information floor remains \(n^{-1}\).

%%% FIELD note-frontier-consumer
The explicit thresholds provide a baseline-only decision rule for comparing another chart-reviewed outcome with another treatment--covariate record when working nuisance models are not trusted.
